% Page budget: 1.35 pages | Owns: gantry coordinates, first-principles equations, scheduling, LPV-LFR realisation, and physical parameterisation.
% WRITING GUIDE
% Purpose: Define the physical coordinates and derive the exact continuous-time LPV-LFR baseline used by every later section.
% Primary reference for Section II-B and its appendix material:
% docs/Thesis-documentation/LPV_LFR_Stepwise_Derivation/main.tex
% (the version shared with Jasper, including the G-Delta interconnection figure).
% Follow its signal definitions, direct mass-equation construction, block partitioning,
% loop-closing verification, and allocation of detailed algebra to the appendix.
% Supporting checks/details only: LPV_LFR_internal_loop_Well_posedness,
% LPV_Mass_invertible, LPV_LFR_Justification_Drenth, LPV_LFR_Derivation,
% and docs/fp-model-structure.md. Where their presentation differs, follow
% LPV_LFR_Stepwise_Derivation for Section II-B.
% Sources: README "Source map per subsection", rows II-A, II-B, II-C and II Frozen-LTI baseline; mandatory papers and the reference gate as in the README.
% Research-plan anchor: Preserve Aspect 1's explicit position dependence; update the original discretisation wording to the final method.
% Current decisions: Continuous-time physics, endogenous scheduler Y, fixed-step RK4, and identifiable parameter combinations.
% Choices to justify: Coordinate frame, Y as scheduler, continuous-time formulation, RK4 step and substeps, full six-channel LFR realization, and reduced physical parameterization. Give the physical or numerical reason for each choice, not only its definition.
% Cautions: The final pipeline uses the full six-channel scheduling loop, not the experimental six-to-four SVD reduction. Resolve the supervisor comment about dual-gantry terminology; do not copy unverified bibliography entries.
% Planned figures: Physical coordinate schematic and compact interconnection between G and Delta(Y).
% Section is finished when: The LFR collapses to the original model, well-posedness assumptions are stated, and notation matches Section 03.
%
% LPV-LFR DERIVATION ROADMAP
% 1. Define q=[X,Theta,Y]^T, the force transformation, measured outputs, assumptions, and the admissible range of Y.
% 2. State M(Y)qdd+Cqd+Kq=f and expose the structural dependence M(Y)=M0+Y M1+Y^2 M2.
% 3. Construct the LFR directly from the polynomial mass equation:
%       a1=Y qdd, a2=Y a1, M0 qdd + M1 a1 + M2 a2 = f_net,
%       z=[qdd; a1], w=[a1; a2] (stepwise latent names; qdd is not renamed to a).
% 4. Collect w=Delta(Y)z with Delta(Y)=Y I6 and give the constant block-level G realization with all signal dimensions.
% 5. Eliminate w and z in a short exactness check to recover the original equation of motion.
% 6. State the well-posedness condition over the complete scheduling range and relate it to invertibility of M(Y); place the longer proof outside the main text.
% 7. Only after the conceptual construction, mention that the implementation
%    evaluates the resulting rational expression in Horner form. Put
%    M(Y)^(-1)=[N0+Y N1+Y^2 N2]/d(Y), the explicit Ni entries, and the
%    determinant coefficients in the appendix or technical supplement.
% 8. Keep this derivation continuous-time; Section III-A owns the RK4 scheme, Section V only its step size.
%
% DERIVATION WARNING
% Do not start from v=M(Y)^(-1)f_net. Follow LPV_LFR_Stepwise_Derivation:
% construct the repeated scheduling loop from qdd, a1=Y qdd, a2=Y a1 with
% z=[qdd; a1], w=[a1; a2], then verify it by closing the loop. Older inverse-first notes
% may be used only as supporting algebra. Verify terminology against the
% primary literature by Drenth and Hoekstra.
% MUST ESTABLISH (II-B after eq:delta, agreed 2026-10-06):
% 1. G is a constant LTI system in Drenth Eq. (6) form; its first block row is the state equation.
% 2. Closing the algebraic loop gives back eq:eom exactly.
% 3. Well posed iff M(Y) nonsingular, true for every Y; hence closed-form M(Y)^{-1}=N(Y)/d(Y)
%    (N_i, d(Y) and Horner evaluation in the appendix). The reduced-parameter bound stays in II-C.
% MUST ESTABLISH (II-C, agreed 2026-10-07):
% 1. The fourteen raw parameters enter M0,M1,M2,C,K only through theta_base (eq:phi); this work's own.
% 2. Argument in the main text, ordered toward the conclusion: input-output behaviour fixes the matrices;
%    twelve entries, three equal m_h, ten independent; so exactly ten combinations are structurally
%    identifiable; only then named theta_base (eq:phi). Model-only (no data), no new symbols, no
%    Proposition/proof environment (agreed 2026-10-07).
% 3. Four lost directions in one sentence; estimates reported in theta_base.
% 4. Positive theta_base does not give M(Y)>0; that holds iff eq:lfr_admissibility (proof in app:lfr-wellposed).
% 5. One sentence: training keeps eq:lfr_admissibility at every iterate. The m_Delta map (eq:mdelta_map)
%    and its justification are stated in Section III-C (agreed 2026-10-07), not here.
% Compact form agreed 2026-10-07: no separate definition paragraph, no null-direction display.
% Moved to Section V as bullets: rho value, combination-error metric with the m_Delta normalisation, raw vs
% reduced coordinates per arm, practical identifiability on the records. Appendix app:params keeps the table only.
% Main text: assumptions, polynomial dependence, direct LFR construction, compact G, exactness, well-posedness statement, and one implementation sentence.
% Appendix/supplement: expanded M_i and N_i matrices, determinant polynomial, longer invertibility/well-posedness proof, and symbolic verification.
% MUST ESTABLISH (II-A, PROPOSED cycle 06, inferred from README ownership and source map row II-A):
% 1. The system is Garcia-Herreros' dual-drive gantry; q=(X,Theta,Y) avoids a closed kinematic chain and
%    exposes the load-distribution coupling (Garcia Sec. 2.2); P links q to the measured actuator frame.
% 2. eq:eom with Garcia's small-angle and no-Coriolis simplifications (Sec. 2.3); Y enters only M(Y).
% 3. Adapted, not adopted: all three coordinates kept (Garcia's (11) keeps X, Theta), Coulomb terms of
%    Eq. 6 omitted with a pointer to Section V (D-204); L_b fixed because it defines q.
% MUST ESTABLISH (II frozen-LTI baseline, PROPOSED cycle 06): the LPV baseline at nominal parameters is
%    named; the frozen-LTI comparator evaluates M at a fixed Y_op (eq:lti), is not trained, and is a
%    comparator only (D-010, D-242 step 1). Y_op value in Section V.
% CYCLE 06 NOTES (header only):
% - Resolved existing todos: "Dynamics not represented..." (owned by Section III opening, not repeated);
%   coordinate reason and "nonuniform ... moving payload" (Garcia Sec. 2.2 cited in prose); "geometric
%   dimensions" (L_b, d defined); "Look at this section" (rewrite); Telica parameter-identification
%   remark (simulation only, THESIS-RESULTS Sec. 2); operating range of Y (M(Y)>0 for every Y, ranges
%   in Section V); "source for self-scheduled" (Toth Def. 7.2); SVD reduction (header Caution and
%   meeting audit: excluded); "Left here" (draft complete). Coulomb insertion todo replaced by the
%   pointer to Section V and a conflict todo (D-022 vs D-204); its Hoekstra Cond. 8 candidate is
%   unverified and dropped.
% - Removed displays eq:lfr_mass_factorization and eq:lfr_chain (no reference elsewhere; the content
%   is in eq:lfr_constant and eq:delta). New labels: sec:system, eq:wellposed, eq:lti.
% - d(Y) renamed det M(Y) (d is the payload offset); appendix symbol a removed (a_1, a_2 are latent signals).
% - Left to other sections: RK4 step and substeps (III-A owns the scheme, V the value); Y_op value,
%   stroke and data ranges, practical identifiability (Section V); m_Delta map (III-C).
% - Stale in other files (not edited): 01_introduction.tex Aspect 1 todo ("LPV versus frozen LTI
%   exists only as a baseline check") is now covered by eq:lti; 06_results.tex line 23 todo idem.
%   99_appendix.tex app:lfr-details edited to follow eq:lfr_G (partition no longer repeated).
\section{Dual-Drive Gantry System and Physics Baseline}
\ifdraft
\begin{itemize}
\item Opening: derives the baseline that Section~III augments; delivers the first contribution of Section~I; names II-A, II-B, II-C (style profile)
\end{itemize}
\fi

This section derives the baseline that Section~\ref{sec:augmentation}
augments: a continuous-time LPV-LFR realisation of the dual-drive gantry with
rational dependence on the payload position, the first contribution of
Section~I.
Section~\ref{sec:system} states the equations of motion,
Section~\ref{sec:lfr} their LPV-LFR realisation and its well-posedness, and
Section~\ref{sec:params} the parameter combinations that input-output data
identify.
\todo{Missing: contribution wording. Section~\ref{sec:params} also delivers the ten identifiable combinations and the admissibility condition, which the Introduction's first contribution does not name. Candidate: ``\dots with rational dependence on payload position, parameterised by its ten identifiable parameter combinations'' (D-191, D-229).}

\subsection{System and first-principles baseline}\label{sec:system}
\ifdraft
\begin{itemize}
\item Two X drives carry a cross-arm through flexible joints; a Y drive moves the payload along it (Garcia Sec.~2.2)
\item $q=(X,\Theta,Y)$ avoids a closed kinematic chain and shows the load-distribution coupling; $P$ links it to the actuator frame (Garcia Sec.~2.2; adopted)
\item Equations of motion with $\cos\Theta\approx1$, $\sin\Theta\approx0$ and no Coriolis term (Garcia Sec.~2.3); $Y$ enters only $M(Y)$
\item Adapted: all three coordinates kept, unlike Garcia's (11); Coulomb terms of Garcia Eq.~(6) omitted (reason in Section~V); $L_b$ fixed because it defines $q$
\item Nominal values from Garcia Table~I, in the appendix
\end{itemize}
\fi
The system is the dual-drive gantry of Garc\'ia-Herreros et
al.~\cite{garcia2013model} (Fig.~\ref{fig:gantry_system}).
Two parallel linear drives carry a cross-arm through flexible joints.
A third drive moves the payload along the cross-arm.
\todo{Missing: the supervisor comment on dual-gantry terminology (header Caution) is not in the decisions or the meeting audit. Candidate: ``dual-drive gantry'', the term of~\cite{garcia2013model} (title); the Introduction still writes ``dual-gantry''.}
\begin{figure}[b]
  \centering
  \def\gantryabsorber{0}% baseline only; the absorber version is in Section III
  \input{figures/gantry_system_with_absorber}
  \caption{Lumped-parameter gantry schematic adapted from the mechanical
  arrangement in \cite{garcia2013model}. Red labels denote actuator coordinates
  and blue labels generalised coordinates.}
  \label{fig:gantry_system}
\end{figure}

We use the generalised coordinates $q=[X\;\Theta\;Y]^\top$ of
Garc\'ia-Herreros et al., which avoid a closed kinematic chain and make the
coupling caused by the non-uniform load distribution
explicit~\cite[Sec.~2.2]{garcia2013model}.
The drives act and are measured in the actuator frame, with forces
$\uact=[F_{X1}\;F_{X2}\;F_Y]^\top$ and positions
$\qact=[X_1\;X_2\;Y]^\top$.
The two frames are related by
\begin{equation}
  q = P^{-\top}\qact,
  \qquad
  u = P\uact,
  \qquad
  P =
  \begin{bmatrix}
    1 & 1 & 0\\
    L_b/2 & -L_b/2 & 0\\
    0 & 0 & 1
  \end{bmatrix},
\label{eq:Ptransform}
\end{equation}
where $u\in\R^3$ is the generalised force and $L_b$ the cross-arm length
between the flexible joints.
Thus $X=(X_1+X_2)/2$ is the cross-arm translation, $\Theta=(X_1-X_2)/L_b$ its
yaw angle, and $Y$ the payload position measured from the centre of the
cross-arm~\cite[Sec.~2.2]{garcia2013model}.

Following Garc\'ia-Herreros et al., we set $\cos\Theta\approx1$ and
$\sin\Theta\approx0$, since $\Theta$ stays within tens of microradians in
operation~\cite[Sec.~2.3]{garcia2013model}.
As they do, we also neglect the Coriolis and centripetal
terms~\cite[Sec.~2.3]{garcia2013model}.
The equations of motion are then
\begin{equation}
  M(Y)\ddot{q}(t) + C\dot{q}(t) + K q(t) = u(t),
  \label{eq:eom}
\end{equation}
with
{\footnotesize
\setlength{\arraycolsep}{3pt}
\begin{equation}
M(Y)=
\begin{bmatrix}
m_1+m_2+m_b+m_h & \tfrac{L_b}{2}(m_1-m_2)-m_hY & 0\\
\tfrac{L_b}{2}(m_1-m_2)-m_hY &
\substack{J_b+J_h+\tfrac{L_b^2}{4}(m_1+m_2)\\{}+m_h(d^2+Y^2)} & -m_hd\\
0 & -m_hd & m_h
\end{bmatrix},
\label{eq:mass_matrix}
\end{equation}
\begin{equation}
\begin{aligned}
C&=
\begin{bmatrix}
c_{g1}+c_{g2} & \tfrac{L_b}{2}(c_{g1}-c_{g2}) & 0\\
\tfrac{L_b}{2}(c_{g1}-c_{g2}) &
c_{b1}+c_{b2}+\tfrac{L_b^2}{4}(c_{g1}+c_{g2}) & 0\\
0 & 0 & c_y
\end{bmatrix},\\[1ex]
K&=
\begin{bmatrix}
0&0&0\\
0&k_{b1}+k_{b2}&0\\
0&0&0
\end{bmatrix},
\end{aligned}
\label{eq:damping_stiffness_matrices}
\end{equation}
}
where $M:\R\to\R^{3\times3}$ is the inertia matrix,
$C,K\in\R^{3\times3}$ are the damping and stiffness matrices, $m_1$, $m_2$,
$m_b$ and $m_h$ the masses of the two X drives, the cross-arm and the payload,
$J_b$ and $J_h$ the yaw inertias of the cross-arm and the payload, $d$ the
offset of the payload from the cross-arm, $c_{g1}$, $c_{g2}$ and $c_y$ the
damping of the drives, $c_{b1}$, $c_{b2}$, $k_{b1}$ and $k_{b2}$ the damping
and stiffness of the flexible joints, and $\theta\in\R^{14}$ collects these
fourteen parameters, with the nominal values
of~\cite[Table~I]{garcia2013model} in Appendix~\ref{app:params}.
The payload position enters only $M(Y)$: it shifts the coupling between $X$
and $\Theta$ and adds $m_hY^2$ to the yaw inertia.

The baseline adapts the model of Garc\'ia-Herreros et al.\ in three choices.
Their simplified model keeps only $X$ and
$\Theta$~\cite[Eq.~(11)]{garcia2013model}.
We keep all three coordinates, because $Y$ schedules $M(Y)$ and the model
predicts all three measured positions.
We omit the Coulomb friction terms
of~\cite[Eq.~(6)]{garcia2013model}, for a data-design reason given in
Section~\ref{sec:setup}.
The length $L_b$ is fixed and not part of $\theta$, because it defines the
coordinates $q$ in \eqref{eq:Ptransform}.
\todo{Missing: source conflict on the Coulomb terms. D-022 keeps the baseline as Garc\'ia-Herreros et al.\ derived it, which includes them; D-204 omits them so that friction stays an effect the augmentation must capture. Candidate: keep the omission with the pointer to Section~\ref{sec:setup} and cite D-204 there (README source map row II-A).}

\subsection{Self-scheduled LPV-LFR representation}\label{sec:lfr}
\ifdraft
\begin{itemize}
\item $Y$ is a state and the scheduling variable (quasi-LPV, T\'oth Def.~7.2); $M(Y)^{-1}$ makes the model rational in $Y$, which an LPV-LFR represents exactly (Drenth thesis Sec.~2.1)
\item $M(Y)=M_0+YM_1+Y^2M_2$; $Y$ enters only through $Y\ddot q$ and $Y^2\ddot q$, realised by repeating $Y$, which keeps the coupling of $Y$ and $Y^2$ (this work; precedent Drenth thesis Eq.~(2.9), Sec.~2.1.1)
\item Latent variables $z,w$ and Delta-block $\Dsch(Y)=YI_6$, the full six channels; not claimed minimal (this work)
\item LTI part $G$ in the continuous-time form of Drenth thesis Eq.~(2.1); its first block row is the state equation (adapted)
\item Closing the algebraic loop gives back \eqref{eq:eom} exactly (this work)
\item Well posed iff $M(Y)$ nonsingular, true for every $Y$; closed-form $M(Y)^{-1}=N(Y)/\det M(Y)$, no numerical solve (Drenth thesis Eq.~(2.4); this work)
\item $\mathcal M_{\mathrm{LPV}}$ is this baseline at the nominal parameters; $\mathcal M_{\mathrm{LTI}}$ freezes $Y$ in $M$, a comparator for Aspect~1, not trained (D-010, D-242)
\end{itemize}
\fi
Since $Y$ is a component of the state, \eqref{eq:eom} is a quasi-LPV model
with scheduling variable $Y$~\cite[Def.~7.2]{toth2010modeling}, which Drenth
et al.\ call self-scheduled~\cite[Sec.~3.2]{drenth2025lpvlfr}.
Solving \eqref{eq:eom} for $\ddot q$ requires $M(Y)^{-1}$, whose entries are
rational in $Y$.
An LPV-LFR represents such rational dependence exactly, as a constant LTI
system $G$ in feedback with a block $\Dsch$ that is linear in the scheduling
variable~\cite[Sec.~2.1]{drenth2025thesis} (Fig.~\ref{fig:lfr}).
We construct this realisation for the gantry directly from the polynomial
structure of $M(Y)$.

\begin{figure}[t]
  \centering
  \input{figures/lfr_loop}
  \caption{Self-scheduled baseline interconnection. The Delta-block closes
  the latent variables of $G$ through $w=\Dsch(Y)z$, with
  $\Dsch(Y)=YI_6$ and $Y$ a component of the state.}
  \label{fig:lfr}
\end{figure}

The inertia matrix \eqref{eq:mass_matrix} is quadratic in $Y$,
\begin{equation}
  M(Y)=M_0+YM_1+Y^2M_2,
  \label{eq:Msplit}
\end{equation}
where $M_0=M(0)$ and $M_1,M_2\in\R^{3\times3}$ depend on $m_h$ only
(Appendix~\ref{app:lfr-details}).
$Y$ therefore enters \eqref{eq:eom} only through $Y\ddot q$ and
$Y^2\ddot q=Y(Y\ddot q)$.

As in Drenth's rational embedding~\cite[Eq.~(2.9)]{drenth2025thesis}, we
realise $Y^2$ by applying $Y$ twice, not as a second scheduling variable.
Independent scheduling variables would admit combinations of $Y$ and $Y^2$
that cannot occur~\cite[Sec.~2.1.1]{drenth2025thesis}.
With the latent signals $a_1=Y\ddot q\in\R^3$ and $a_2=Ya_1\in\R^3$,
\eqref{eq:eom} becomes
\begin{equation}
  M_0\ddot q+M_1a_1+M_2a_2=u-C\dot q-Kq,
  \label{eq:lfr_constant}
\end{equation}
which contains no $Y$.
The scheduling is collected in the Delta-block
\begin{equation}
  w=\begin{bmatrix}a_1\\a_2\end{bmatrix}
   =\underbrace{\begin{bmatrix}YI_3&0\\0&YI_3\end{bmatrix}}_{\Dsch(Y)=YI_6}
    \begin{bmatrix}\ddot q\\a_1\end{bmatrix}
   =\Dsch(Y)z,
  \label{eq:delta}
\end{equation}
where $z,w\in\R^6$ are the latent variables and
$\Dsch:\R\to\R^{6\times6}$.
We do not claim this realisation to be minimal.

Solving \eqref{eq:lfr_constant} for $\ddot q$ with the constant $M_0^{-1}$
gives the LTI part in the continuous-time form
of~\cite[Eq.~(2.1)]{drenth2025thesis},
\begin{equation}
  \begin{bmatrix}\dot{\tilde x}\\ z\\ q\end{bmatrix}
  =\underbrace{\begin{bmatrix}
  A_x&B_w&B_u\\ C_z&D_{zw}&D_{zu}\\ C_y&0&0
 \end{bmatrix}}_{G}
  \begin{bmatrix}\tilde x\\ w\\ u\end{bmatrix},
  \label{eq:lfr_G}
\end{equation}
where $\tilde x=\col(q,\dot q)\in\R^6$ is the state, $q$ the output and
$G\in\R^{15\times15}$ is constant (blocks in
Appendix~\ref{app:lfr-details}).
Its first block row is the state equation.

Closing the loop with \eqref{eq:delta} gives $a_1=Y\ddot q$ and
$a_2=Y^2\ddot q$.
Substituted in \eqref{eq:lfr_constant}, these return \eqref{eq:eom}, so the
LFR realises the baseline exactly.
The LFR is well posed if $I_6-D_{zw}\Dsch(Y)$ is nonsingular for every
$Y$~\cite[Eq.~(2.4)]{drenth2025thesis}.
For our realisation, we show in Appendix~\ref{app:lfr-details} that
\begin{equation}
  \det\bigl(I_6-D_{zw}\Dsch(Y)\bigr)=\frac{\det M(Y)}{\det M_0},
  \label{eq:wellposed}
\end{equation}
so the LFR is well posed exactly where $M(Y)$ is nonsingular.

The baseline is well posed for every payload position, because positive
masses and inertias give $M(Y)\succ0$ for every $Y\in\R$
(Appendix~\ref{app:lfr-wellposed}).
The implementation therefore solves the loop in closed form,
$\ddot q=M(Y)^{-1}(u-C\dot q-Kq)$, with
$M(Y)^{-1}=N(Y)/\det M(Y)$ and both polynomials quadratic in $Y$
(Appendix~\ref{app:lfr-details}).

We call the baseline \eqref{eq:delta}, \eqref{eq:lfr_G} at the nominal
parameters $\mathcal M_{\mathrm{LPV}}$.
Section~\ref{sec:aug_structure} discretises this continuous-time model.
To measure what the scheduling contributes (Aspect~1 of Section~I), we compare it with a frozen LTI
baseline $\mathcal M_{\mathrm{LTI}}$, which evaluates the inertia at a fixed
payload position $Y_{\mathrm{op}}\in\R$,
\begin{equation}
  M(Y_{\mathrm{op}})\ddot q(t)+C\dot q(t)+Kq(t)=u(t),
  \label{eq:lti}
\end{equation}
with $Y_{\mathrm{op}}$ given in Section~\ref{sec:setup}.
$\mathcal M_{\mathrm{LTI}}$ is a comparator only: it is neither trained nor
augmented.
\todo{Missing: model notation. Candidate: $\mathcal M_{\mathrm{LPV}}$ and $\mathcal M_{\mathrm{LTI}}$, since $M$ is the inertia matrix (README open conflict on model names; Section~\ref{sec:aug_closed_loop} uses $\mathcal M_{\mathrm U}$).}
\todo{Missing: no runner evaluates $\mathcal M_{\mathrm{LTI}}$ yet; the frozen-$Y$ option of the baseline block is a capability only (README source map, II Frozen-LTI baseline). Candidate: add the step~1 evaluation of THESIS-RESULTS before Section~\ref{sec:results} reports it.}

\subsection{Physical parameters and identifiable combinations}\label{sec:params}
\ifdraft
\begin{itemize}
\item The input-output behaviour fixes $M_0,M_1,M_2,C,K$ and conversely, so the identifiable part of $\theta$ is what the matrices contain (this work)
\item Twelve nonzero entries, three equal to $m_h$, leave ten independent functions: exactly ten combinations are identifiable (this work; definition Ovchinnikov et al.\ Def.~7)
\item Written physically as $\theta_{\mathrm{base}}$; four lost directions; estimates reported in $\theta_{\mathrm{base}}$ (this work; Gautier and Khalil counterpart unverified (?))
\item Positive $\theta_{\mathrm{base}}$ does not give $M(Y)\succ0$; that holds iff \eqref{eq:lfr_admissibility} (this work; proof Appendix~\ref{app:lfr-wellposed})
\item Training keeps \eqref{eq:lfr_admissibility} at every iterate; the map and its justification are in Section~\ref{sec:aug_closed_loop}
\end{itemize}
\fi
The fourteen parameters $\theta$ enter the model only through $M_0$, $M_1$,
$M_2$, $C$ and $K$.
Two parameter vectors with the same input-output behaviour satisfy
\eqref{eq:eom} along the same trajectories.
The gantry is fully actuated, so every triple $(q,\dot q,\ddot q)$ occurs on
some trajectory.
Hence they give the same matrices.
Conversely, equal matrices give the same model.
The identifiable part of $\theta$ is therefore what these matrices contain.

Up to symmetry, \eqref{eq:mass_matrix} and \eqref{eq:damping_stiffness_matrices}
have twelve nonzero entries.
Three of them, $M_{0,33}$, $-M_{1,12}$ and $M_{2,22}$, equal $m_h$, which
leaves ten distinct functions of $\theta$.
Their Jacobian with respect to $\theta$ has rank ten.
Exactly ten independent functions of $\theta$ are therefore
identifiable~\cite[Def.~7]{ovchinnikov2021computing}.
We report estimates in these ten combinations, written in physically
readable form as
\begin{equation}
  \theta_{\mathrm{base}}=[k_{b,\Sigma},c_{g1},c_{g2},c_y,c_{b,\Sigma},
  m_h,m_\Sigma,m_\Delta,J_{\mathrm{eff}},d]^\top\in\R^{10},
  \label{eq:phi}
\end{equation}
with $k_{b,\Sigma}=k_{b1}+k_{b2}$, $c_{b,\Sigma}=c_{b1}+c_{b2}$,
$m_\Sigma=m_1+m_2+m_b$, $m_\Delta=m_1-m_2$ and
$J_{\mathrm{eff}}=J_b+J_h+\tfrac{L_b^2}{4}(m_1+m_2)$.
The four lost directions are $k_{b1}-k_{b2}$, $c_{b1}-c_{b2}$, $J_b-J_h$, and
a transfer of mass from $m_1$ and $m_2$ to $m_b$ with a compensating
$J_b+J_h$.
\todo{Missing: whether to name $\theta_{\mathrm{base}}$ the gantry counterpart of the base inertial parameters of robotics (Gautier and Khalil 1990, key \texttt{gautier1990minimum}). Candidate: add the clause and the citation after reading the paper; no local PDF, unverified (README source map row II-C).}
\todo{Missing: notation. The index $\Delta$ in $m_\Delta$ clashes with the scheduling block $\Dsch$, and $m_\Delta$ is used in Sections~\ref{sec:aug_closed_loop} and~\ref{sec:obc}. Candidate: $m_{\mathrm{diff}}$, the code name (Dirk decides, all three sections change together).}

Unlike positive masses and inertias (Section~\ref{sec:lfr}), positive
entries of $\theta_{\mathrm{base}}$ do not imply $M(Y)\succ0$.
For $m_h,m_\Sigma,J_{\mathrm{eff}}>0$, $M(Y)\succ0$ for every $Y\in\R$ if
and only if
\begin{equation}
  m_\Sigma J_{\mathrm{eff}}>\tfrac{L_b^2}{4}\,m_\Delta^2
  \label{eq:lfr_admissibility}
\end{equation}
(Appendix~\ref{app:lfr-wellposed}).

An unconstrained estimate of $\theta_{\mathrm{base}}$ can violate
\eqref{eq:lfr_admissibility}.
Section~\ref{sec:aug_closed_loop} therefore trains $\theta_{\mathrm{base}}$
in coordinates that keep \eqref{eq:lfr_admissibility}, and with it
well-posedness, at every iterate.
